\begin{description}
  \item[Appendix A: Post's lattice cheat sheet; completeness tables.] A compact reference for the clone-theoretic classification of Boolean function classes, including the principal nodes and inclusions in Post's lattice, standard closure properties, and completeness criteria used throughout the book. This appendix is intended as a quick lookup table for functional completeness arguments, separation results, and recurring examples. It is organized for consultation rather than exposition, and it should be read alongside the chapters on closure, clones, and Boolean function classes.

  \item[Appendix B: Operad/PROP identities; coherence diagrams.] A reference sheet for the algebraic identities and coherence laws that recur in the main text, including operadic substitution, PROP composition, symmetry and naturality constraints, and diagrammatic equalities. The emphasis is on the identities needed to verify that composite constructions are well-typed and that alternative parenthesizations or wiring orders agree. This appendix collects the equations and commuting-diagram patterns that support routine coherence checks in the book's compositional arguments.

  \item[Appendix C: ZX, traced monoidal, decorated cospan guides.] A compact guide to the diagrammatic formalisms used in the book: ZX-calculus conventions, traced monoidal feedback notation, and decorated cospan interfaces for open systems. This appendix serves as a bridge between the abstract categorical language and the string-diagram or network representations used in examples and case studies. It is intended as a notation guide and a reminder of the standard translation patterns between algebraic and graphical presentations.

  \item[Appendix D: Metrics catalog for \(\varepsilon(r)\): defs, units, calibration.] A catalog of the error and robustness metrics used to quantify approximation quality across resolutions, including the definition of \(\varepsilon(r)\), the associated units and normalization conventions, and the calibration procedures used when comparing artifacts at different scales. The intent is to make metric usage consistent across chapters and to support reproducible reporting of numerical results. Where multiple conventions are possible, this appendix records the book's chosen convention and the conditions under which comparisons are meaningful.
\end{description}

These appendices are supplementary reference material. They are designed to be consulted alongside the main chapters rather than read sequentially, and they collect the compact statements, identities, and conventions that support the book's closure, compositionality, and symmetry arguments.

\medskip
\noindent\textbf{Use and scope.} The appendices are intentionally concise: they prioritize lookup value over exposition, and they summarize the conventions, identities, and calibration rules that recur across the book. Where a chapter develops a concept in detail, the corresponding appendix entry should be read as a reminder of notation, a checklist of hypotheses, or a compact statement of the relevant law.

\medskip
\noindent\textbf{Cross-reference policy.} Appendix A supports the Boolean closure and completeness arguments in the foundations chapters; Appendix B supports the algebraic and coherence calculations used in operadic and monoidal constructions; Appendix C supports the diagrammatic reasoning used in feedback, reversible, and quantum-flavored sections; and Appendix D supports the quantitative comparisons used in robustness, evaluation, and reproducibility discussions. Together, these appendices provide the reference layer for the book's formal and computational claims.

\medskip
\noindent\textbf{Reading note.} If a result in the main text depends on a convention, identity, or metric normalization, the corresponding appendix should be treated as the authoritative local reference. In particular, the appendices are meant to stabilize notation across chapters and to reduce duplication of standard material in the main narrative.

\medskip
\noindent\textbf{Scope of the appendix set.} Appendix A is the lookup table for Boolean function classes and completeness criteria; Appendix B is the identity sheet for operads, PROPs, and coherence; Appendix C is the notation guide for ZX-calculus, traced monoidal structure, and decorated cospans; and Appendix D is the metric reference for \(\varepsilon(r)\), including units, normalization, and calibration. Together they form the book's compact reference layer for formal statements, diagrammatic conventions, and quantitative comparisons.

\medskip
\noindent\textbf{Consultation guide.} When using these appendices, prefer the narrowest entry that matches the claim at hand: use Appendix A for Boolean-function classification and completeness checks, Appendix B for algebraic identity verification, Appendix C for diagrammatic translations and feedback conventions, and Appendix D for scale-dependent error reporting. This keeps the main chapters uncluttered while preserving a single authoritative location for each recurring convention.
